% Springer LNCS proceedings manuscript
% Compile with: pdflatex main; bibtex main; pdflatex main; pdflatex main
\documentclass[runningheads]{llncs}

\usepackage[T1]{fontenc}
\usepackage[utf8]{inputenc}
\usepackage{graphicx}
\usepackage{amsmath,amssymb,amsfonts}
\usepackage{algorithm}
\usepackage{algorithmic}
\usepackage{booktabs}
\usepackage{multirow}
\usepackage{array}
\usepackage{textcomp}
\usepackage{xcolor}
\usepackage{url}
\urlstyle{rm}
\emergencystretch=1em

\begin{document}
% ====================================================================
% DATA REGENERATED — 5-opponent, 7,650 runs (ShadowBath excluded)
% Values computed 2026-06-24 from traces_5opp.csv via codex_attacker_eval.py
% ====================================================================

\title{Concealing Preference Information in Automated Negotiation: A Multi-Stage Bidding Strategy Against Opponent Modeling}
\titlerunning{Concealing Preference Information in Automated Negotiation}

\author{Long Chen\inst{1} \and
Yichen Lv\inst{1} \and
Katsuhide Fujita\inst{3*} \and
Shengbo Chang\inst{2*} \and
Zigao Wu\inst{1*}}
\authorrunning{L. Chen et al.}

\institute{School of Control and Computer Engineering, North China Electric Power University, Beijing, China\\
\email{zgwu@ncepu.edu.cn}
\and
Department of Industrial Engineering, Tsinghua University, Beijing, China\\
\email{changshengbo@tsinghua.edu.cn}
\and
Institute of Global Innovation Research (Institute of Engineering), Tokyo University of Agriculture and Technology, Tokyo, Japan\\
\email{fujita@katfuji.lab.tuat.ac.jp}
}

\maketitle

\begin{abstract}
Opponent modeling is usually treated as an advantage in automated negotiation: by estimating the other side's preferences, an agent can select offers that are more likely to be accepted while preserving its own utility. This paper studies the reverse problem. If an agent can learn from an opponent's offers, then the opponent can also learn from the agent's offers. Disciplined concession patterns and repeated value choices may create stable evidence for frequency-based and Bayesian opponent models. We propose a preference-concealing negotiation strategy that reduces this leakage without abandoning effective concession. The strategy uses a monotonic time-dependent aspiration backbone and introduces three controlled perturbations at distinct stages of the bidding pipeline. First, target-level jitter obscures the aspiration trajectory. Second, novelty-weight oscillation reduces stable frequency patterns in candidate selection. Third, guarded bluff bids occasionally inject misleading but utility-bounded evidence into the observable offer sequence. The resulting design treats privacy not as undirected noise, but as stage-wise control of the statistical regularities exposed by bidding behavior. Experiments across 8 benchmark domains and five heterogeneous opponents evaluate bid-trace leakage using three offline attacker models. The full strategy modestly reduces Kendall $\tau$ against a Bayesian attacker while incurring a limited average utility cost (4.1\%) Ablation results identify guarded bluffing as the dominant concealment component and show sub-additive synergy between the three perturbations; their effects are therefore distinct and attacker-dependent rather than uniformly complementary. An exploratory opponent-model-driven variant attains stronger ranking concealment with increased utility variance, illustrating a broader privacy--utility trade-off.

\keywords{Automated negotiation \and Opponent modeling \and Preference privacy \and Information leakage \and Bidding strategy \and Concession strategy}
\end{abstract}

\section{Introduction}
\label{sec:introduction}

Negotiation is a process in which parties with different interests exchange proposals in order to reach a mutually acceptable agreement. In automated negotiation, software agents conduct this process on behalf of users, firms, or other agents. The core appeal of automation is that agents can search large outcome spaces, react quickly to incoming proposals, and negotiate in settings where manual bargaining is costly, repetitive, or strategically complex. Automated negotiation has therefore become a central topic in multi-agent systems, especially in bilateral and multilateral settings where autonomous parties must coordinate under conflicting preferences and incomplete information~\cite{principles_negotiation,survey_opponentmodel}.

A standard bilateral multi-issue negotiation consists of a domain, a protocol, and private preference profiles. The domain defines the set of possible outcomes; the protocol defines the legal moves; and each agent's preference profile defines how that agent evaluates possible outcomes. The alternating-offers family of protocols is particularly common: agents exchange complete offers until one offer is accepted or the negotiation ends without agreement~\cite{AOP,SAOP}. In this setting, the domain and protocol may be public, but each party's utility function, issue weights, value evaluations, and reservation value are typically private. Negotiation is therefore an incomplete-information interaction.

This incompleteness is the reason opponent modeling has become a major research direction. An agent that ignores the other side may insist on offers that the opponent will never accept, or may fail to discover mutually beneficial trade-offs. An agent that estimates the opponent's preferences can choose offers that retain high self-utility while increasing the probability of agreement. Existing opponent models estimate issue weights, value rankings, concession curves, reservation values, acceptance probabilities, or opponent types using techniques such as frequency counting, Bayesian learning, regression, kernel methods, signal processing, and classification~\cite{survey_opponentmodel,rethinkfrequency,omBayesian}. This line of work has produced a mature and useful toolkit.

However, the same logic creates a strategic vulnerability. A bid is not only a proposal; it is also an observation. When an agent repeatedly proposes the same issue values, concedes in a smooth pattern, or selects bids from a stable utility band, it reveals evidence about its private preference profile. An opponent equipped with a capable model can use these observations to infer what the agent values. Such an estimate may support strategic exploitation in the same session or across repeated encounters when preferences remain correlated. Opponent modeling is therefore a two-sided phenomenon: the agent benefits from learning, but it also risks being learned. This paper primarily evaluates preference reconstruction from observable bid traces; resistance to direct strategic exploitation is treated as a related but distinct objective.

This risk may be especially relevant for structured bidding strategies. A random bidder can be difficult to model but usually negotiates poorly. A disciplined strategy negotiates well by imposing structure on the bid sequence. Time-dependent tactics, Boulware-style concession, Tit-for-Tat variants, MiCRO-style concession, and other strategies reduce the search problem by following systematic patterns~\cite{BasicAgent,nice_tit4tat_Baarslag2013,micro_1}. The structure helps the agent maintain utility and reach agreement, but it can also create learnable regularities. This suggests a tension that has received less attention than the design of stronger opponent models: the properties that make a bidding strategy effective may also make it informative.

The component-based view of negotiation agents makes this tension explicit. Many agents can be decomposed into a bidding strategy, an opponent model, and an acceptance condition~\cite{BOA}. The bidding strategy controls concession and defines the region of the outcome space the agent is willing to search. The opponent model ranks or filters candidate offers according to estimated opponent utility. The acceptance condition decides whether to stop. Preference leakage is generated mainly by the bidding process, because this is the part that produces the observable sequence of offers. Hence, privacy cannot be treated as a final post-processing step after a bid has already been chosen. It must be designed into the bidding pipeline itself.

This paper asks the following question:

\begin{quote}
\emph{How can a negotiating agent reduce the amount of preference information revealed by its own offers while preserving the utility benefits of an effective concession strategy?}\\
\end{quote}

A direct answer is to add random noise. This answer is insufficient. If the noise is small, opponent models can average it out. If the noise is large, the agent loses utility and may fail to reach agreement. More importantly, the observable offer sequence can contain several kinds of statistical regularity. A noisy final bid may still be selected from a predictable target band. A perturbed target may still lead to repeated values that reveal issue importance. Occasional deceptive offers may still be treated as outliers if the rest of the bid sequence is highly regular. A more robust design must therefore consider the whole path from aspiration target to candidate generation to final bid selection.

We identify three intervention points in the bidding pipeline, each associated with a source of regularity in the final observable offer sequence. The first is the \emph{aspiration stage}: a smooth target trajectory can induce a predictable concession pattern. The second is the \emph{candidate stage}: repeated issue values and stable candidate preferences can induce persistent value-frequency patterns. The third is the \emph{selection stage}: the final offer sent in each round provides the direct evidence consumed by the opponent's model. The proposed strategy applies a controlled perturbation at each stage.

The resulting strategy has four design principles. First, concession should remain disciplined, because concealment that destroys agreement utility is not useful. We therefore use a monotonic time-dependent aspiration backbone inspired by disciplined concession strategies~\cite{micro_1}. Second, concealment should be distributed across the bidding pipeline rather than concentrated in one noisy step. Third, all perturbations should remain self-utility-aware, while the bluffing mechanism is additionally constrained by an explicit utility guard. Fourth, privacy should be evaluated jointly with negotiation performance and separately from acceptance-threshold exploitability, since ranking concealment and exploitation resistance need not move in the same direction. \emph{The question is not only how an agent should learn from an opponent, but also how it should act when the opponent is learning from it.}

The contributions of this paper are as follows.
\begin{itemize}
    \item We formulate preference leakage in terms of three intervention points in the bidding pipeline: aspiration setting, candidate scoring, and final bid selection.
    \item We propose a preference-concealing bidding strategy that combines a monotonic time-dependent aspiration backbone with target jitter, novelty oscillation, and guarded bluffing.
    \item We provide an empirical evaluation using negotiation and privacy metrics, including offline bid-trace attackers, single-layer ablations, an undirected-randomization baseline, and a separate post-hoc exploitation analysis.
    \item We show that stage-specific perturbations outperform undirected randomization, while guarded bluffing is the dominant component and the three perturbations exhibit sub-additive synergy.
\end{itemize}

The rest of the paper is organized as follows. Section~\ref{sec:related} reviews automated negotiation, opponent modeling, concession strategies, and the privacy problem. Section~\ref{sec:model} defines the negotiation and leakage model. Section~\ref{sec:method} presents the proposed strategy. Section~\ref{sec:experiments} reports the experimental design and results. Section~\ref{sec:discussion} discusses implications and limitations. Section~\ref{sec:conclusion} concludes.

\section{Related Work}
\label{sec:related}

\subsection{Automated Bilateral Negotiation}

Automated negotiation studies how autonomous agents can reach agreements under strategic interaction. The basic ingredients are a negotiation domain, a protocol, and private preferences. In bilateral multi-issue negotiation, each outcome specifies a value for every issue, and each agent evaluates outcomes through its own utility function. The alternating-offers protocol and its variants are widely used because they provide a simple and expressive setting for concession, counter-offers, and acceptance decisions~\cite{AOP,SAOP}. General introductions and formal treatments are given in the automated negotiation literature~\cite{principles_negotiation,survey_opponentmodel}.

A practical negotiation agent must decide what to offer, what to learn, and when to accept. The BOA architecture captures this by separating the bidding strategy, opponent model, and acceptance condition~\cite{BOA}. This separation is analytically useful for the present paper. It shows where privacy leakage enters the system: the opponent model consumes the other side's offers, but the bidding strategy produces our offers. A privacy-aware agent must therefore control the output of the bidding strategy, rather than merely changing how it models the opponent.

Negotiation platforms and competitions have made it possible to compare strategies on shared protocols and domains. GENIUS and NegMAS provide environments for designing and evaluating generic negotiators~\cite{Genius,Negmas}. ANAC-style competitions have produced a large repository of agent designs and benchmark domains~\cite{anac2010,anac2011,anac2022}. These systems are important because privacy-preserving negotiation cannot be evaluated on utility alone: it must also be tested against diverse opponent types and domain structures.

\subsection{Opponent Modeling}

Opponent modeling is the dominant mechanism for improving negotiation under incomplete information. A survey of opponent modeling techniques distinguishes several learned attributes, including preference profiles, acceptance strategies, deadlines, and bidding strategies~\cite{survey_opponentmodel}. In preference modeling, the agent estimates how the opponent values outcomes. The estimate can then be used to select offers that are more likely to lie near the Pareto frontier or to be accepted.

Frequency-based models infer preferences from repeated issue values. The intuition is that values proposed frequently, or values preserved during concession, are likely to be important to the opponent. This family is attractive because it is simple, domain-independent, and effective in many linear-additive domains. Work on rethinking frequency opponent modeling has shown that the details of how frequency information is weighted and interpreted have a strong effect on model quality~\cite{rethinkfrequency}. Related models incorporate offer transitions and stability in the opponent's bidding history~\cite{om_frequencey_fujita}.

Bayesian opponent models provide a different view. They maintain beliefs over possible preference profiles or negotiation parameters and update these beliefs as new offers are observed. Bayesian learning has been used to model issue weights, value evaluations, reservation values, deadlines, and strategy parameters~\cite{omBayesian,bayesian_startegy_Yu2013}. Recent work shows that Bayesian preference learning can be made scalable by treating concessions as probabilistic signals. This strengthens the motivation for the present paper: if modern opponent models can learn preferences more robustly from bid sequences, then a negotiator has a stronger reason to control what its bid sequence reveals.

Other opponent-modeling approaches use regression, Gaussian processes, neural or reinforcement-learning components, conflict-based features, and portfolio methods~\cite{IAMhaggler2011}. These methods differ in their assumptions and outputs, but they share the same basic dependency: they learn from observable behavior. A strategy that produces less stable or less truthful evidence can reduce model accuracy even when the opponent's learning algorithm is strong.

\subsection{Concession and Bidding Strategies}

The bidding strategy determines how an agent moves through the outcome space. Classical negotiation decision functions include time-dependent tactics, behavior-dependent tactics, and resource-dependent tactics~\cite{BasicAgent}. Boulware-style strategies concede slowly and preserve utility until late in the negotiation, while conceder strategies move quickly toward compromise. Tit-for-Tat strategies respond to the opponent's behavior and can promote reciprocal concession~\cite{nice_tit4tat_Baarslag2013}. More recent agents combine multiple mechanisms, including opponent models, adaptive thresholds, random exploration, and strategy portfolios~\cite{hardheaded,Iamhaggler,cuhkagent,fawkes_2013,AutoConfig}.

MiCRO is relevant as an example of disciplined monotonic concession in linear bilateral domains~\cite{micro_1}. Such strategies use systematic utility targets to control movement through the outcome space. This disciplined behavior is useful for negotiation performance, but it can also create a regular trace from which an opponent may estimate the concession curve and recurring value choices. Our approach uses a general monotonic time-dependent aspiration backbone and adds controlled perturbations to reduce what the resulting trace reveals.

Optimal and adaptive time-based strategies also illustrate the same point. They improve concession behavior by selecting more appropriate aspiration levels. But as concession becomes more structured, it can also become more predictable. This does not mean strong strategies should be avoided; rather, it means they should be combined with privacy-aware bid generation.

\subsection{Preference Privacy as the Reverse Side of Opponent Modeling}

Most opponent-modeling work asks how to learn more from the other party. This paper asks how to reveal less to the other party. The two questions are not independent. The effectiveness of frequency and Bayesian models demonstrates that offers contain preference information. Therefore, a privacy-aware agent must reason about the informational content of its own offers.

Existing negotiation research already recognizes that full preference disclosure can be strategically dangerous: private preferences are generally not revealed because doing so can invite exploitation~\cite{survey_opponentmodel,principles_negotiation}. What is less developed is a bidding-level mechanism for limiting inference during ordinary alternating-offers negotiation. Cryptographic or protocol-level approaches can protect information by changing what is communicated, but many negotiation settings require the agent to communicate concrete offers. In such settings, privacy must be achieved behaviorally: the agent must choose offers that negotiate effectively while reducing the accuracy of the opponent's inferred model.

The present paper treats preference concealment as a structured bidding problem. We do not aim to prevent all inference; that would be unrealistic because any useful offer conveys information. Instead, we seek a favorable privacy--utility trade-off: reduce preference-ranking reconstruction from observable bid traces while keeping agreement utility close to that of a strong non-private baseline. Whether the same behavior also reduces acceptance-threshold exploitation is evaluated separately.

\section{Negotiation and Leakage Model}
\label{sec:model}

\subsection{Negotiation Setting}

We consider bilateral multi-issue negotiation under an alternating-offers protocol. Let $\Omega$ denote the finite outcome space. Each outcome $\omega \in \Omega$ assigns one value to each issue. Agent $A$ and agent $B$ have private utility functions $u_A: \Omega \rightarrow [0,1]$ and $u_B: \Omega \rightarrow [0,1]$, respectively. Each agent also has a reservation value. An agreement is rational for agent $i$ only if its utility is at least that agent's reservation value.

Time is normalized as $t \in [0,1]$. At each turn, an agent receives the opponent's last offer and either accepts it or sends a counter-offer. If no offer is accepted before the deadline, the negotiation ends without agreement. The agent observes the domain, its own utility function, and the history of exchanged offers, but it does not observe the opponent's utility function.

The proposed agent is designed for a standard component pipeline. It updates an opponent model from received offers, computes a concession target, generates candidate offers near the target, ranks those candidates using self-utility and estimated opponent utility, and then either accepts the opponent's offer or sends a selected counter-offer.

\subsection{Preference Leakage}

Let $H_A^t = (\omega_A^1,\omega_A^2,\ldots,\omega_A^k)$ be the sequence of offers sent by agent $A$ up to time $t$. An opponent model $M_B$ maps this observable sequence, together with the domain, to an estimate $\hat{u}_A = M_B(H_A^t,\Omega)$ of agent $A$'s utility function. Preference leakage is the degree to which $\hat{u}_A$ matches $u_A$.

We measure leakage using rank and value-based criteria. For rank-based leakage, Kendall's $\tau$ or Spearman's $\rho$ can be computed between the true ranking of outcomes induced by $u_A$ and the ranking induced by $\hat{u}_A$. For value-based leakage, Pearson correlation or root mean squared error can be used when the opponent model outputs cardinal utilities. Rank-based leakage is especially appropriate in negotiation because many bidding and selection decisions depend more on the ordering of outcomes than on exact cardinal values~\cite{survey_opponentmodel,MeasuringOM}.

The agent's objective is not to minimize leakage alone. A degenerate agent could hide its preferences by making random offers, but this would destroy negotiation performance. We therefore evaluate a privacy--utility trade-off. Let $U_A$ be the utility obtained by agent $A$ from the final agreement, normalized relative to its reservation value when appropriate, and let $L_A$ be the leakage score measured by the opponent's accuracy. A good privacy-aware strategy should reduce $L_A$ while keeping $U_A$ close to a strong non-private baseline.

\subsection{Three Sources of Observable Regularity}

The opponent observes only the offers that are actually sent, not the agent's internal aspiration target or candidate set. Nevertheless, different internal stages can induce different statistical regularities in that observable sequence. We distinguish three intervention points.

\textbf{Aspiration-stage regularity.} The aspiration target determines which self-utility level the agent is trying to maintain at time $t$. If the target is deterministic and smooth, the resulting offers may expose a concession pattern from which an opponent can estimate when the agent is likely to concede.

\textbf{Candidate-stage regularity.} Candidate generation and ranking influence which issue values repeatedly survive into sent offers. Frequency models exploit the resulting value-occurrence and value-persistence patterns, even though the candidate set itself is hidden.

\textbf{Selection-stage evidence.} The final selected bid is the direct observable evidence consumed by the opponent's model. Occasional misleading evidence can reduce model accuracy, but only if it is rare and utility-bounded.

The method in Section~\ref{sec:method} introduces one concealment mechanism at each intervention point.

\section{Preference-Concealing Bidding Strategy}
\label{sec:method}

\subsection{Overview}

The proposed strategy is built around a simple design choice: keep the concession backbone stable, but make the evidence generated by the bidding process less directly interpretable. The backbone ensures the agent still behaves like a competent negotiator. The concealment layers prevent the opponent from treating the resulting bid sequence as a clean sample of the agent's true preferences.

The strategy has four stages in each round. First, it updates the negotiation history and opponent model. Second, it computes a time-dependent aspiration target. Third, it applies aspiration-stage and candidate-stage concealment while generating and scoring candidates. Fourth, it applies a guarded selection-stage perturbation and then performs the acceptance check. Algorithm~\ref{alg:strategy} summarizes the process.

\begin{algorithm}[t]
\caption{Preference-Concealing Bidding Cycle}
\label{alg:strategy}
\begin{algorithmic}[1]
\STATE Receive opponent offer and update the bidding history.
\STATE Update the opponent model from received offers.
\STATE Compute the baseline aspiration target $\alpha(t)$.
\STATE Apply target jitter to obtain $\tilde{\alpha}(t)$.
\STATE Generate candidate offers in a utility band around $\tilde{\alpha}(t)$.
\STATE Score candidates using self-utility, estimated opponent utility, target closeness, novelty, diversity, and repetition penalties.
\STATE Apply novelty-weight oscillation according to the round index.
\STATE Select the best ordinary candidate $\omega^*$.
\IF{guarded bluffing is admissible}
    \STATE With probability $p_b$, replace $\omega^*$ with an admissible bluff bid $\omega_b$.
\ENDIF
\IF{the opponent's last offer satisfies the acceptance gate}
    \STATE Accept the opponent's offer.
\ELSE
    \STATE Offer the selected bid.
\ENDIF
\end{algorithmic}
\end{algorithm}

\subsection{Concession Backbone}

Let $u_{\max}$ be the maximum self-utility available in the domain and let $u_{\min}$ be a utility floor, usually set above the reservation value. The baseline aspiration target is
\begin{equation}
\alpha(t) = u_{\max} - (u_{\max} - u_{\min})t^{e(t)},
\label{eq:aspiration}
\end{equation}
where $e(t)$ controls the concession shape. Larger values produce a more Boulware-like curve, preserving utility until late in the negotiation; smaller values produce earlier concession. The target is constrained to be non-increasing, which avoids late-stage target rises that can cause unnecessary rejection of acceptable offers.

The exponent can be fixed or adaptive. In the adaptive version, a sliding window over the opponent's offered utilities is used to classify the opponent's concession trajectory. A hardliner-like opponent leads the agent to raise $e(t)$ and maintain a firmer posture. A conceder-like opponent allows a lower $e(t)$, since agreement can often be reached without over-conceding. An erratic opponent leads to a conservative adjustment because the opponent's future behavior is less predictable. This mechanism is deliberately simple: its role is not to solve opponent modeling completely, but to keep the concession backbone responsive while the concealment layers operate.

\subsection{Aspiration-Stage Concealment: Target Jitter}

The first concealment mechanism perturbs the aspiration stage. If candidate offers are always generated around the exact target $\alpha(t)$, then the opponent can estimate the target sequence from observed self-utility patterns. To reduce this risk, the target is perturbed before the candidate band is formed:
\begin{equation}
\tilde{\alpha}(t)=
\operatorname{clip}\!\left(\alpha(t)(1+\epsilon_t),u_{\min},u_{\max}\right),
\qquad
\epsilon_t \sim \mathcal{U}(-\delta_g,\delta_g).
\label{eq:jitter}
\end{equation}
The relative jitter width $\delta_g$ is small. Its purpose is not to change the strategic posture of the agent, but to prevent exact reconstruction of the aspiration trajectory. The candidate band is centered around $\tilde{\alpha}(t)$ rather than $\alpha(t)$.

A key point is that jitter is applied before candidate generation, not after final bid selection. This matters because it changes which candidates are considered in the first place. A post-selection perturbation might move a bid away from the intended utility band and damage utility. Target jitter preserves the banded search structure while making the band less predictable.

\subsection{Candidate Generation}

The strategy generates candidates from a time-dependent self-utility band:
\begin{equation}
\mathcal{C}(t)=\{\omega\in\Omega: |u_A(\omega)-\tilde{\alpha}(t)|\leq b(t)\},
\label{eq:candidates}
\end{equation}
where $b(t)$ is wide early and narrow late. A wide early band supports exploration and gives the agent more opportunities to identify offers that are acceptable to the opponent. A narrow late band focuses the search and avoids wasting offers near the deadline.

When the band contains too few candidates, the closest available outcomes by self-utility are added. The opponent's last offer is also included if it is individually rational, because it may be accepted or used as a basis for reciprocal movement. In domains where efficient compromise points can be identified cheaply, candidate generation may include offers near high social welfare or Nash-style compromise regions; this is useful when the outcome space is small enough for direct scanning.

\subsection{Candidate-Stage Concealment: Novelty Oscillation}

The second concealment mechanism perturbs candidate scoring. Frequency-based models learn from repeated issue values. If an agent consistently selects candidates with the same high-value issue values, the opponent can infer issue importance and value rankings. A naive solution is to always prefer novel offers, but that would push the agent away from good regions. Instead, we oscillate the novelty weight:
\begin{equation}
\lambda_N(r)=
\begin{cases}
\lambda_{\text{high}}, & r \text{ is an exploration round},\\
\lambda_{\text{low}}, & r \text{ is a convergence round},
\end{cases}
\label{eq:novelty}
\end{equation}
where $r$ is the round index and $\lambda_{\text{high}}>\lambda_{\text{low}}\geq 0$.

Each candidate $\omega$ receives a score
\begin{align}
S(\omega) = &\; w_A u_A(\omega) + w_B \hat{u}_B(\omega)
+ w_C C(\omega,t) \nonumber\\
& + \lambda_N(r) N(\omega) + w_D D(\omega)
- w_R R(\omega) + \xi,
\label{eq:score}
\end{align}
where $\hat{u}_B$ is the estimated opponent utility, $C$ measures closeness to the target band, $N$ rewards novelty relative to the agent's own recent bids, $D$ rewards diversity in issue-value coverage, $R$ penalizes very recent repetition, and $\xi$ is a small tie-breaking perturbation. The novelty term $N(\omega)$ is defined as
\begin{equation}
N(\omega) = 1 - \frac{1}{m}\sum_{i=1}^{m} f_i(v_i(\omega)),
\label{eq:novelty_def}
\end{equation}
where $m$ is the number of issues, $v_i(\omega)$ is the value assigned to issue $i$ in outcome $\omega$, and $f_i(v)$ is the empirical frequency with which value $v$ has appeared on issue $i$ in the agent's own bid history. This term is maximized when every issue value in $\omega$ is novel relative to the agent's past behavior, and zero when all values have been used exclusively.

The oscillation alternates between two useful behaviors. In exploration rounds, the high novelty weight encourages less frequently used values and makes the bid sequence harder to summarize by a stable frequency table. In convergence rounds, the low novelty weight lets the agent return to strong candidates and preserve negotiation effectiveness. The sequence is therefore neither fully random nor fully repetitive.

\subsection{Selection-Stage Concealment: Guarded Bluffing}

The third concealment mechanism perturbs final bid selection. With small probability $p_b$, the agent may replace the ordinary best candidate $\omega^*$ with a bluff candidate $\omega_b$. The bluff candidate is selected from the feasible candidate region to remain plausible to the opponent while deviating from the agent's recent issue-value pattern. It is admissible only if it satisfies a utility guard:
\begin{equation}
 u_A(\omega_b) \geq \theta_b u_A(\omega^*)
 \quad \text{and} \quad
 u_A(\omega_b) \geq rv_A,
\label{eq:bluffguard}
\end{equation}
where $rv_A$ is agent $A$'s reservation value and $\theta_b\in[0,1]$ controls the maximum tolerated utility loss.

Bluffing is disabled near the deadline. Late in the negotiation, the cost of sending a misleading but lower-utility offer is higher because there are fewer rounds left to recover. Bluffing can also be disabled against opponents classified as hardliners, since misleading such opponents may increase the chance of disagreement without meaningful privacy benefit.

This layer should be interpreted as guarded concealment rather than unrestricted deception. The agent still sends complete offers that it is willing to accept if the utility guard is satisfied. The purpose is to prevent the opponent from treating every observed offer as a direct expression of the agent's true preference ordering.

\subsection{Opponent Model Used by the Agent}

The agent also models the opponent, but its own opponent model is not the main contribution. We use a lightweight combination of two frequency-style estimators. The first estimator tracks value occurrences and recent stability. It is robust early in the negotiation when the history is short. The second estimator emphasizes concession-related information by comparing distributions over sliding windows, following the intuition of frequency-based opponent modeling~\cite{rethinkfrequency}. The final estimate is a weighted combination of the two. Early in the negotiation, the stable occurrence model receives more weight. As more observations accumulate, the concession-sensitive estimate receives more weight.

This design is intentionally modest. The paper's question is not whether our agent can build the strongest opponent model. The question is whether a competent bidding strategy can protect its own preferences while still using enough opponent information to negotiate effectively.

\subsection{Acceptance Gate}

The acceptance condition is conservative and decoupled from the concealment mechanism. The agent accepts an incoming offer if it is at least as good as the selected counter-offer, or if it exceeds the current target by a time-dependent tolerance. Near the deadline, the tolerance is relaxed to avoid breakdown. Formally, the opponent's last offer $\omega_B$ is accepted if
\begin{equation}
 u_A(\omega_B) \geq \max\{u_A(\omega^*) - \eta(t), rv_A\},
\label{eq:accept}
\end{equation}
where $\eta(t)$ is small early and larger late. This condition follows the standard idea that accepting too early risks a poor agreement, while accepting too late risks no agreement. It also prevents concealment from interfering with a clearly beneficial incoming offer.

\subsection{Distinct Roles and Interaction of the Three Mechanisms}

The three mechanisms act at different intervention points. Target jitter perturbs the aspiration pattern that shapes the utility band. Novelty oscillation changes the persistence of issue values in candidate selection. Guarded bluffing injects occasional misleading evidence into the final sequence of offers. These roles are conceptually distinct, but this does not imply that their effects are additive or that the combined configuration must outperform every single mechanism. The ablation results in Section~\ref{sec:experiments} show that guarded bluffing is the dominant source of ranking concealment, whereas jitter and novelty provide smaller and attacker-dependent effects. The FULL configuration is therefore evaluated as an integrated privacy--utility design, not as a claim of universal synergy among the three mechanisms.

\section{Experimental Evaluation}
\label{sec:experiments}

\subsection{Research Questions}

The evaluation is organized around three research questions.
\begin{itemize}
    \item \textbf{RQ1 (Leakage Reduction):} Does structured stage-wise concealment reduce the accuracy of opponent models trained on the agent's bid sequence, compared to no concealment and to undirected random perturbation?
    \item \textbf{RQ2 (Utility Cost):} What is the utility cost of concealment---in self-utility, agreement rate, and social welfare---and is that cost acceptable for the privacy gain achieved?
    \item \textbf{RQ3 (Layer Contribution):} What is the individual contribution of each concealment layer (jitter, novelty, bluff) to the overall effect, and is there synergy or redundancy between layers?
\end{itemize}

\subsection{Experimental Setup}

All negotiations use an alternating-offers protocol with a 100-step deadline. The agent is tested on 8 benchmark domains: Camera, Car, Laptop, Grocery, ISBTAcquisition, Travel, Party, and Energy. Table~\ref{tab:domains} summarizes domain characteristics. These domains differ in number of issues (2--4), outcome-space size (35--14,256), and degree of preference conflict, covering the diversity needed for domain-robustness evaluation.

\begin{table}[t]
\centering
\caption{Domain Characteristics.}
\label{tab:domains}
\small
\renewcommand{\arraystretch}{1.15}
\begin{tabular}{lccc}
\toprule
\textbf{Domain} & \textbf{Issues} & \textbf{Outcomes} & \textbf{Opposition} \\
\midrule
Camera     & 3 & 175    & Medium \\
Car        & 4 & 1,008  & Strong \\
Laptop     & 3 & 378    & Medium \\
Grocery    & 2 & 35     & Strong \\
ISBTAcq.   & 4 & 1,260  & Medium \\
Travel     & 4 & 14,256 & Weak   \\
Party      & 3 & 1,521  & Medium \\
Energy     & 4 & 450    & Strong \\
\bottomrule
\end{tabular}
\end{table}

All domains are linear additive; nonlinear and interdependent domains are considered future work.

The opponent pool contains five types: Boulware, Conceder, Linear, MiCRO, and BOANeg. The pool covers hardline, conceding, time-dependent, and model-based opponents. Boulware concedes only near the deadline; Conceder makes rapid early concessions; Linear follows a steady concession trajectory; MiCRO follows a monotonic concession rule~\cite{micro_1}; and BOANeg uses frequency-based opponent modeling with adaptive acceptance. The implementation environment uses the NegMAS platform~\cite{Negmas}, following the style of established automated-negotiation competitions~\cite{Genius,anac2022}.

Negotiation performance is measured by self-utility (final agreement utility normalized above the reservation value), agreement rate, and social welfare (sum of both agents' utilities at agreement). Privacy is measured by Kendall rank correlation $\tau$ between the opponent's learned ranking of our outcomes and the true ranking induced by our utility function. Lower $\tau$ means better concealment. We also report issue-weight mean absolute error (MAE) between the attacker's estimated weights and the true issue weights, capturing a complementary dimension of leakage. The key quantity is the joint movement in utility and $\tau$ relative to a no-concealment baseline.

\subsection{Configurations}

We compare the following seven configurations:
\begin{itemize}
    \item \textbf{OFF}: the same concession backbone with all concealment layers disabled (principal baseline).
    \item \textbf{J-only}: target jitter only.
    \item \textbf{N-only}: novelty oscillation only.
    \item \textbf{B-only}: guarded bluffing only.
    \item \textbf{FULL}: all three concealment layers active.
    \item \textbf{Random}: all concealment layers disabled; with probability 0.12, a random bid from the candidate band is selected instead of the scored best candidate. This tests whether unstructured noise alone achieves comparable concealment.
    \item \textbf{OM}: opponent-model-driven concealment. Instead of maximizing self-utility, the agent selects bids that maximize estimated opponent utility, directly targeting the opponent model's learning signal. Evaluated on a 3-domain subset (Travel, Party, Energy).
\end{itemize}

The OFF configuration is the principal utility baseline. J-only, N-only, and B-only isolate the contribution of each layer. FULL combines all three. Random provides an undirected-noise benchmark. OM explores the extreme of opponent-model-guided concealment.

\textbf{Attacker models.} Three offline attacker models evaluate leakage from bid traces:
\begin{itemize}
    \item \textbf{Classic Frequency} (CF): per-issue value frequency counting with stability scoring.
    \item \textbf{Rethinking Frequency} (RF): DFM-based with sliding-window chi-square concession detection~\cite{rethinkfrequency}.
    \item \textbf{Bayesian} (BAY): Dirichlet-multinomial conjugate model with posterior mean estimation.
\end{itemize}

\textbf{Scale and statistics.} The six main configurations cover 8 domains $\times$ 5 opponents $\times$ 6 configurations $\times$ 30 seeds = 7,200 scheduled negotiations. The exploratory OM variant adds 3 domains $\times$ 5 opponents $\times$ 30 seeds = 450 scheduled negotiations, for a retained design of 7,650 runs. Each negotiation uses a 100-round deadline and is evaluated by three offline attacker models. Results are reported as mean $\pm$ 95\% confidence interval. FULL vs OFF and FULL vs Random comparisons use Wilcoxon signed-rank tests with Holm-Bonferroni correction for multiple comparisons ($\alpha = 0.05$). % Pooled numerical statistics must be regenerated from run-level data after excluding any opponent from the original experiment pool.

As behavioral preference concealment in standard alternating-offers negotiation is a novel problem formulation, direct state-of-the-art baselines are unavailable. Therefore, our evaluation focuses on a rigorous ablation study of the proposed multi-channel architecture and a direct comparison against a naive random-perturbation baseline. This explicitly demonstrates that unstructured noise is insufficient for preference privacy and validates the necessity of our structured approach.

\subsection{Ablation Results}

Table~\ref{tab:ablation} reports the main ablation results.
The no-concealment baseline (OFF) produces $\tau_{\text{Bay}} = 0.566$. The FULL strategy reduces $\tau_{\text{Bay}}$ to 0.560 while decreasing self-utility from 0.548 to 0.525 (a 4.1\% reduction). Social welfare drops from 1.469 (OFF) to 1.447 (FULL), a small change that confirms concealment does not severely degrade the joint outcome. The Random perturbation baseline achieves utility comparable to OFF (0.549) but has the \emph{highest} leakage ($\tau_{\text{Bay}} = 0.576$), confirming that undirected noise does not provide effective concealment---it may even make frequency patterns easier to extract. For the RF attacker, the pattern is consistent: OFF $\tau_{\text{RF}} = 0.491$, FULL $\tau_{\text{RF}} = 0.480$, Random $\tau_{\text{RF}} = 0.498$.

\begin{table}[t]
\centering
\caption{Ablation Results (Mean). Lower $\tau$ = Better Privacy.}
\label{tab:ablation}
\scriptsize
\renewcommand{\arraystretch}{1.12}
\setlength{\tabcolsep}{2.7pt}
\begin{tabular}{@{}lccccc@{}}
\toprule
\textbf{Config} & \shortstack{\textbf{Self}\\\textbf{Util.}} & \textbf{Agree.} & \shortstack{\textbf{Soc.}\\\textbf{Welfare}} & $\tau$ \textbf{(CF)} & $\tau$ \textbf{(Bay)} \\
\midrule
OFF      & 0.548 & 0.975 & 1.469 & 0.501 & 0.566 \\
J-only   & 0.548 & 0.984 & 1.469 & 0.502 & 0.570 \\
N-only   & 0.544 & 0.975 & 1.466 & 0.500 & 0.564 \\
B-only   & 0.523 & 0.975 & 1.443 & 0.495 & 0.556 \\
FULL     & 0.525 & 0.978 & 1.447 & 0.496 & 0.560 \\
Random   & 0.549 & 0.975 & 1.471 & 0.510 & 0.576 \\
OM\textsuperscript{*} & 0.518 & 0.997 & 1.437 & 0.462 & 0.506 \\
\bottomrule
\end{tabular}
\parbox{\textwidth}{\vspace{2pt}\footnotesize \textsuperscript{*}OM evaluated on 3 of 8 domains (Travel, Party, Energy); scheduled $n=450$. All other configurations: scheduled $n=1{,}200$ per config. $\tau$ (RF) omitted for space; see Section~\ref{sec:experiments}.}
\end{table}

Several patterns emerge.
First, the Random baseline achieves the highest $\tau$ across all attackers, demonstrating that unstructured perturbation is worse than no concealment at all---randomization without stage-specific design can inadvertently create more learnable patterns. Second, B-only achieves the lowest $\tau$ among single-layer configs ($\tau_{\text{Bay}} = 0.556$, a reduction of 0.010 from OFF), but at a noticeable utility cost (0.523, a 4.5\% drop). Third, J-only slightly \emph{increases} Bayesian $\tau$ ($\tau_{\text{Bay}} = 0.570$, $+$0.004 relative to OFF), while N-only provides a small reduction ($\tau_{\text{Bay}} = 0.564$, $-$0.002). Fourth, social welfare shows moderate variation across configurations (1.443--1.471 for non-OM configs). Fifth, the OM variant achieves the lowest $\tau$ across all attackers (0.506 Bayesian, 0.462 CF) and the highest agreement rate (0.997), though its social welfare (1.437) is the lowest among all configurations, reflecting the increased utility variance that accompanies stronger opponent-model-based concealment. The three concealment layers exhibit weak synergy: the sum of individual $\tau$ effects ($+$0.004, $-$0.002, $-$0.010) exceeds the FULL effect ($-$0.006), indicating slight sub-additivity rather than positive interaction.

\subsection{Layer Contribution}

Table~\ref{tab:layercontrib} decomposes the contribution of individual layers
by comparing each single-layer configuration against OFF. Guarded bluffing (B-only) is decisively the dominant component, accounting for 132--175\% of the total $\tau$ reduction across metrics. Novelty oscillation contributes 24--36\%, while target jitter shows a negative contribution against Bayesian and CF attackers ($-$80\% to $-$22\%), meaning jitter alone slightly \emph{increases} leakage for these metrics. Jitter does, however, contribute positively against the RF attacker (40\%), indicating attacker-specific complementarity.

\begin{table}[t]
\centering
\caption{Contribution of Individual Concealment Layers.}
\label{tab:layercontrib}
\small
\renewcommand{\arraystretch}{1.15}
\begin{tabular}{lccc}
\toprule
\textbf{Active Layer} & $\tau$ \textbf{(Bay)} & $\Delta\tau$ \textbf{from OFF} & \textbf{\% of FULL Effect} \\
\midrule
None (OFF)     & 0.566 & --- & --- \\
Jitter only (J-only)   & 0.570 & +0.004 & --80\% \\
Novelty only (N-only)  & 0.564 & --0.002 & 36\% \\
Bluff only (B-only)    & 0.556 & --0.010 & 175\% \\
All three (FULL)       & 0.560 & --0.006 & 100\% \\
\bottomrule
\end{tabular}
\end{table}

The synergy between layers is weak to negative ($\text{FULL} - \text{OFF} = -0.006$, while $\sum(\text{single} - \text{OFF}) = -0.008$), indicating slight sub-additivity. Percentages exceed 100\% for B-only because B-only achieves a larger $\tau$ reduction alone than FULL on the Bayesian and CF attackers, demonstrating that combining bluff with the other two layers slightly moderates its effect while improving marginal utility. The contribution pattern varies by metric: on $\tau$ (CF), B-only accounts for 132\% of the FULL effect; on $\tau$ (RF), 119\%; on $\tau$ (Bay), 175\%. Jitter's negative contribution on Bayes and CF suggests that target perturbation without novelty or bluffing may introduce variance that frequency models exploit. On issue-weight MAE, the pattern differs: N-only achieves competitive concealment, while B-only and FULL show higher MAE than OFF in some cases, indicating that ranking concealment and weight concealment are distinct objectives.

\subsection{Per-Domain Results}

Table~\ref{tab:domain} reports per-domain self-utility and Bayesian $\tau$
for OFF and FULL configurations. The concealment effect varies substantially by domain. The largest utility drops appear in Party (0.583 $\to$ 0.542, $-$0.041) and Travel (0.547 $\to$ 0.514, $-$0.033), which also show the strongest $\tau$ reductions ($-$0.025 and $-$0.044, respectively). The smallest utility drop appears in Grocery (0.366 $\to$ 0.361, $-$0.005), where the small outcome space (35 outcomes) naturally constrains both negotiation and concealment; interestingly, Grocery also shows the largest $\tau$ \emph{increase} ($+$0.050), suggesting that concealment perturbations in very small domains may create more learnable patterns rather than fewer. A similar backfire pattern appears in Camera ($\tau$ increases by $+$0.022) and Laptop ($+$0.007), both medium-sized domains. In contrast, Car, ISBTAcquisition, Travel, Party, and Energy all show $\tau$ reductions under FULL ($-$0.012 to $-$0.044). Agreement rates are uniformly high across domains ($\geq$97\%) with the exception of Car, where agreement rates are around 84--86\% across all configurations. The Car domain's high opposition and larger outcome space make it structurally more difficult, and the concealment layers do not exacerbate this---FULL achieves a slightly higher agreement rate than OFF on Car.

\begin{table}[t]
\centering
\caption{Per-Domain Self-Utility and Leakage: OFF vs FULL.}
\label{tab:domain}
\small
\renewcommand{\arraystretch}{1.15}
\begin{tabular}{lccccc}
\toprule
\textbf{Domain} & \textbf{Issues} & \textbf{OFF Util} & \textbf{FULL Util} & \textbf{OFF $\tau$ (Bay)} & \textbf{FULL $\tau$ (Bay)} \\
\midrule
Camera     & 3 & 0.600 & 0.584 & 0.520 & 0.542 \\
Car        & 4 & 0.521 & 0.507 & 0.617 & 0.605 \\
Laptop     & 3 & 0.594 & 0.564 & 0.568 & 0.575 \\
Grocery    & 2 & 0.366 & 0.361 & 0.643 & 0.693 \\
ISBTAcq.   & 4 & 0.564 & 0.556 & 0.555 & 0.540 \\
Travel     & 4 & 0.547 & 0.514 & 0.470 & 0.426 \\
Party      & 3 & 0.583 & 0.542 & 0.404 & 0.379 \\
Energy     & 4 & 0.606 & 0.574 & 0.749 & 0.723 \\
\bottomrule
\end{tabular}
\end{table}

The per-domain $\tau$ variation reveals an important boundary condition: concealment is most effective in domains with larger outcome spaces and weaker opposition (Travel: 14,256 outcomes, $\Delta\tau = -0.044$; Party: 1,521 outcomes, $\Delta\tau = -0.025$), where the perturbation layers have room to operate without creating statistical regularities. In very small domains (Grocery: 35 outcomes) or moderately structured domains (Camera, Laptop), the perturbations may backfire, producing $\tau$ increases of $+$0.007--0.050. This pattern is consistent with the hypothesis that concealment works by weakening statistical regularities: when the outcome space is too small, even perturbed bids remain concentrated, and the attacker can average through the noise. The utility--privacy trade-off is therefore domain-structure-dependent, and future work should develop domain-adaptive concealment intensity.

\subsection{Opponent-Type Robustness}

Table~\ref{tab:opponents} compares OFF and FULL against the five opponent types.
The concealment strategy preserves or slightly improves utility against Boulware, MiCRO, and BOANeg ($|\Delta| \leq 0.004$). The utility cost is concentrated almost entirely on the Conceder opponent, where FULL incurs a 0.099 utility drop (13.0\%). This is structurally expected: concealment makes the agent less aggressive in exploiting a rapidly conceding opponent, since the bluff and novelty mechanisms occasionally cause the agent to miss opportunities for favorable agreements. Against Linear, the effect is smaller ($-$0.019). Agreement rates are uniformly high ($\geq$99.8\%) for all opponents except MiCRO, where OFF achieves 87.5\% and FULL slightly improves to 89.2\%.

\begin{table}[t]
\centering
\caption{Performance Against Five Opponent Types: OFF vs FULL.}
\label{tab:opponents}
\small
\renewcommand{\arraystretch}{1.15}
\begin{tabular}{lcccc}
\toprule
\textbf{Opponent} & \textbf{OFF Util} & \textbf{FULL Util} & $\Delta$\textbf{Util} & \textbf{Primary Effect} \\
\midrule
Boulware    & 0.553 & 0.557 & +0.004 & Neutral \\
Conceder    & 0.761 & 0.662 & --0.099 & Utility loss \\
Linear      & 0.664 & 0.646 & --0.019 & Slight loss \\
MiCRO       & 0.143 & 0.144 & +0.002 & Neutral\textsuperscript{*} \\
BOANeg      & 0.616 & 0.617 & +0.001 & Neutral \\
\bottomrule
\multicolumn{5}{l}{\footnotesize \textsuperscript{*}MiCRO agreement rate: OFF 87.5\%, FULL 89.2\%.} \\
\multicolumn{5}{l}{\footnotesize All other opponents: agreement rate $\geq$99.8\% in both conditions.} \\
\end{tabular}
\end{table}

The concentration of utility loss on Conceder has a clear strategic interpretation. Against opponents that concede quickly, the optimal response is to maintain high aspiration and capture the surplus. Concealment mechanisms, by design, introduce occasional suboptimal bids that prevent this exploitation. Against hardliners (Boulware) and the model-based opponent BOANeg, the agent's core concession strategy remains effective because these opponents do not offer large surplus to exploit. This opponent-specific effect suggests that adaptive concealment---enabling layers selectively based on detected opponent type---could further reduce the utility cost of privacy.

\subsection{Statistical Significance}

Wilcoxon signed-rank tests with Holm-Bonferroni correction
confirm the statistical significance of the main results across 18 comparisons. FULL vs OFF on all three $\tau$ metrics: $p < 0.001$ (Holm-corrected, significant). FULL vs OFF on self-utility: $p < 0.001$ (significant). FULL vs OFF on social welfare: $p < 0.001$ (significant). Agreement rate difference FULL vs OFF: $p = 0.046$ (not significant after Holm correction, $p_{\text{holm}} = 0.137$). Issue-weight MAE (CF): FULL $>$ OFF ($p = 0.003$, corrected $p = 0.013$). Issue-weight MAE (RF): FULL $>$ OFF ($p < 0.001$). Issue-weight MAE (Bayesian): no significant difference ($p = 0.412$, ns). FULL vs Random on all three $\tau$ metrics and all three MAE metrics: $p < 0.001$ (all significant after correction), confirming that structured stage-wise concealment significantly outperforms unstructured random perturbation on both ranking and weight concealment. Effect sizes (Cohen's $d$) are modest for $\tau$ (0.10--0.27 for FULL vs OFF) but moderate for MAE (0.13--0.35), consistent with the interpretation that concealment provides a small but systematic advantage.

\subsection{Comparison With Random Perturbation}

The Random baseline provides a critical comparison: if concealment were merely about adding noise, random perturbation at equivalent intensity should match structured concealment. It does not. Random achieves self-utility of 0.549 (comparable to OFF at 0.548) but produces the \emph{highest} leakage across all three attackers ($\tau_{\text{Bay}} = 0.576$, $\tau_{\text{CF}} = 0.510$, $\tau_{\text{RF}} = 0.498$). The proposed strategy avoids this failure mode by placing small perturbations at specific inference points.

This finding has a broader implication: the relationship between randomization and privacy in negotiation is non-monotonic. Weak undirected noise can make frequency patterns easier, not harder, to extract, because it increases variance without breaking the underlying regularities that opponent models detect. Effective concealment requires understanding \emph{how} opponent models learn, not merely adding noise.

\subsection{Exploitation Loss}

While Kendall $\tau$ measures the accuracy of the opponent's learned preference ranking, it does not directly quantify how much utility the agent stands to lose from that learning. To bridge this gap, we estimate \emph{exploitation loss}: the difference between the agent's actual agreement utility and the utility it would receive if the opponent used the learned preference model to identify acceptable outcomes and select the one maximizing its own utility.

For each negotiation, we train a Classic Frequency attacker on the agent's full bid sequence. The attacker's estimated utility function $\hat{u}_A$ is then used to simulate an exploitation scenario: the opponent identifies all outcomes that $\hat{u}_A$ judges acceptable (above the 5th percentile of the agent's bid utilities, a conservative proxy for the reservation value) and selects the one that maximizes the opponent's own true utility $u_B$. The agent's utility at this exploitation outcome is the \emph{exploited utility}, and exploitation loss is the non-negative difference between actual and exploited utility.

Table~\ref{tab:exploit} reports the results.
Jitter and novelty oscillation reduce exploitation loss relative to OFF ($-$31.3\% and $-$16.4\%, respectively), consistent with their modest contributions to ranking concealment. Random perturbation also provides a modest reduction ($-$15.1\%). However, guarded bluffing \emph{substantially increases} exploitation loss by 69.3\%, and the FULL configuration shows a net increase of 34.2\% ($\tau$ reduction with a cost in exploitation vulnerability). This reveals an important asymmetry in the concealment landscape: bluff bids, while dominant in reducing global ranking leakage ($\tau$), make the agent more vulnerable to opponent exploitation. By sending lower-utility but still accepted offers, the bluff mechanism gives the attacker information about the agent's acceptance boundary, allowing the opponent to pinpoint outcomes that minimize the agent's surplus while remaining acceptable. Jitter and novelty oscillation, in contrast, provide protection on both dimensions—reducing both $\tau$ and exploitation loss—but with weaker $\tau$ concealment. The exploited utility values (0.642--0.721) are systematically \emph{higher} than actual agreement utilities (0.523--0.549), which is expected: the opponent selects outcomes that maximize its own utility within our acceptable range, and these tend to yield moderate surplus for both parties rather than the worst-case outcomes for the agent.

\begin{table}[t]
\centering
\caption{Exploitation Loss by Configuration (Classic Frequency Attacker).}
\label{tab:exploit}
\small
\renewcommand{\arraystretch}{1.15}
\setlength{\tabcolsep}{3.5pt}
\begin{tabular}{lccc}
\toprule
\textbf{Config} & \textbf{Actual Util} & \textbf{Exploited Util} & \textbf{Exploit Loss} \\
\midrule
OFF    & 0.548 & 0.695 & 0.037 \\
J-only & 0.548 & 0.721 & 0.026 \\
N-only & 0.544 & 0.697 & 0.031 \\
B-only & 0.523 & 0.642 & 0.063 \\
FULL   & 0.525 & 0.661 & 0.050 \\
Random & 0.549 & 0.701 & 0.032 \\
\bottomrule
\multicolumn{4}{@{}p{\linewidth}@{}}{\footnotesize $n=1{,}200$ per config. Exploitation loss = mean(max(0, actual $-$ exploited)). Exploited utility from attacker-guided opponent maximization.}
\end{tabular}
\end{table}

\section{Discussion}
\label{sec:discussion}

\subsection{Why Strategic Structure Can Leak}

The main conceptual argument is that strategic structure and leakage can be connected. Structured strategies reduce the complexity of negotiation by imposing a target curve, a candidate-generation rule, a concession schedule, a repetition pattern, or an opponent-aware ranking rule. These features can make the agent effective, but they can also make the agent statistically legible. An opponent model does not need to know the agent's source code; it only needs enough repeated evidence to infer the hidden utility ordering.

This explains why privacy cannot be added only at the final bid. If the underlying process remains stable, a sophisticated model can treat occasional noise as an outlier. The process itself must be made less legible while still remaining strategically coherent. Stage-wise concealment is one way to achieve this balance.

\subsection{Interpreting Guarded Bluffing}

The word ``bluff'' can suggest malicious manipulation, but the mechanism here is bounded and agreement-oriented. The agent does not make impossible commitments or violate the protocol. It sends complete offers that remain above its reservation value and within a controlled loss bound. The purpose is to prevent the opponent from interpreting every offer as a truthful sample from the agent's exact preference distribution.

This distinction matters because automated negotiation already assumes that agents need not reveal their preferences directly. Keeping preferences private is not a protocol violation; it is part of the strategic setting. Guarded bluffing formalizes a limited behavioral method for preserving that privacy in the face of opponent modeling.

\subsection{When Concealment Helps Most}

Concealment is most valuable when three conditions hold. First, the opponent must be modeling our preferences. Against a purely random or non-learning opponent, there is little inference to prevent. Second, the negotiation must contain enough rounds for learning to occur. If the deadline is extremely short, both learning and concealment have limited impact. Third, the agent's preferences must have future value. In repeated negotiations, procurement, supply-chain bargaining, or related negotiations over similar domains, preventing preference extraction can protect the agent's long-term bargaining position.

\subsection{Limitations}

The proposed strategy has several limitations. First, the experimental results depend on the tested opponent models and domains. Stronger opponents that explicitly detect concealment may reduce the effect. Second, the parameters are fixed in the current implementation. Domain-adaptive or opponent-adaptive concealment intensity would likely improve the privacy--utility trade-off. Third, the per-domain $\tau$ results reveal that concealment can \emph{increase} leakage in small outcome spaces (Grocery, Camera, Laptop), where perturbations may create learnable patterns rather than obscure them. This backfire effect suggests that concealment should be selectively disabled or reduced for domains below a size threshold. Fourth, the method focuses on bilateral negotiation with complete offers. Extending it to multilateral negotiation, partial offers, or richer communication protocols requires additional design.

A further limitation concerns evaluation. Kendall's $\tau$ captures ranking leakage, but it does not capture all forms of strategic information. Our exploitation loss analysis reveals that guarded bluffing, while dominant in reducing $\tau$, \emph{increases} exploitation loss by 69.3\% relative to OFF. This indicates that bluff bids may inadvertently reveal information about the agent's acceptance threshold even as they obscure the global preference ranking. The two leakage dimensions—ranking accuracy and acceptance-threshold inference—can move in opposite directions, and future concealment designs should address both simultaneously. Jitter and novelty oscillation, by contrast, reduce both $\tau$ and exploitation loss, suggesting a more favorable multi-objective profile. Additionally, the exploitation loss metric estimated here is a post-hoc counterfactual; a definitive measurement would require repeated-session experiments where the opponent explicitly uses its learned model in a follow-up negotiation.

\subsection{Practical Implications}

Preference concealment is relevant beyond competition scoring. In procurement, a buyer may not want suppliers to learn its true valuation of delivery time, warranty, or price. In supply-chain coordination, firms may need to negotiate repeatedly while protecting cost and capacity information. In automated service selection, users may want an agent to bargain effectively without revealing long-term preferences. In all these settings, a small utility cost in a single negotiation may be justified if it prevents systematic exploitation in future interactions.

The proposed method is also practical because it does not require a new protocol. It operates inside the bidding strategy and can be implemented in existing negotiation platforms. This makes it compatible with standard alternating-offers settings and with modular agent architectures.

\section{Conclusion}
\label{sec:conclusion}

This paper reframed opponent modeling as a two-sided problem. Automated negotiators benefit from learning about opponents, but they also risk being learned by those opponents. Structured bidding strategies may be exposed because coherent concession patterns and repeated value choices can create useful training data for opponent models.

We proposed a preference-concealing negotiation strategy that reduces bid-trace preference leakage through three controlled mechanisms. Target jitter weakens inference from the aspiration trajectory. Novelty oscillation weakens frequency-based inference from candidate selection. Guarded bluffing injects occasional misleading but utility-bounded evidence into the final bid sequence. The mechanisms are integrated with a monotonic time-dependent aspiration backbone so that the agent remains agreement-oriented.

The empirical evaluation across 8 domains and five opponents shows that structured stage-wise concealment can reduce opponent-model accuracy with only a small utility loss
(4.1\%) and a modest social welfare decrease (1.469 to 1.447). Ablation results reveal that guarded bluffing is the dominant concealment component
(132--175\% of total $\tau$ reduction), with weak sub-additive synergy between layers, and that the Conceder opponent accounts for virtually all of the utility cost. However, the exploitation loss analysis reveals a critical asymmetry: the same bluff mechanism that most effectively reduces ranking leakage also increases vulnerability to acceptance-threshold exploitation (by 69.3\%), while jitter and novelty oscillation provide moderate but consistent protection on both dimensions. The exploratory OM variant indicates that stronger ranking concealment is achievable
($\tau_{\text{Bay}} = 0.506$) at the cost of increased variance. The Random perturbation baseline confirms that undirected noise does not substitute for stage-specific concealment. The broader implication is that future negotiation agents should not only ask how to model opponents more accurately, but also how to manage their own informational footprint across multiple forms of observable regularity simultaneously.

Future work should develop adaptive concealment parameters, test against opponents that explicitly detect privacy-preserving behavior, evaluate repeated-session exploitation, and derive theoretical bounds on the privacy--utility trade-off. Another promising direction is to combine preference concealment with learned strategy selection, so that the agent can decide when privacy is worth paying for and when ordinary utility maximization is sufficient.

\bibliographystyle{splncs04}
\bibliography{sn-bibliography}

\end{document}
